\DocumentMetadata{
  pdfversion=2.0,
  lang=es-MX,
  pdfstandard=ua-2
}

\documentclass{sener2025}

\title{Prueba de Correcciones Aplicadas}
\subtitle{Verificación de Tablas Cortas y Listas}
\author{Test}
\date{\today}
\institucion{Secretaría de Energía}
\unidad{Unidad de Pruebas}
\setDocumentoCorto{TestCorrecciones}

\begin{document}

\section{Prueba de Tablas Cortas}

Esta tabla corta debe tener encabezados con fondo dorado, texto gris alineado a la izquierda:

\begin{tabladoradoCorto}
  \caption{Tabla de prueba con estilo corregido}
  \label{tab:prueba_corregida}
  \begin{tabular}{H{3cm}G{2cm}G{2cm}G{2cm}}
    \toprule
    \rowcolor{gobmxDorado} \encabezadodorado{Concepto} & \encabezadodorado{Valor 1} & \encabezadodorado{Valor 2} & \encabezadodorado{Nota} \\
    \midrule
    \textbf{Dato A} & 100 & 200 & Ejemplo \\
    \textbf{Dato B} & 150 & 250 & Prueba \\
    \textbf{Dato C} & 200 & 300 & Test \\
    \bottomrule
  \end{tabular}
\end{tabladoradoCorto}

\section{Prueba de Listas Antes de Figura}

Ahora una lista que debe mantenerse compacta:

\begin{itemize}
  \item Primera viñeta de prueba
  \item Segunda viñeta de prueba
\end{itemize}

Esta figura debe aparecer en la siguiente página sin estirar las viñetas anteriores:

\begin{figure}[H]
  {\centering
  \rule{10cm}{8cm}
  \par}
  \raggedright
  \caption{Figura de prueba que salta de página}
  \label{fig:prueba_salto}
\end{figure}

\section{Comparación con Tabla Larga}

Para comparar, aquí hay una tabla larga que debe mantener su estilo original:

\begin{tabladoradoLargo}
  \begin{longtable}{B{3cm}p{2.5cm}p{2.5cm}p{2.5cm}}
    \caption{Tabla larga de comparación}\label{tab:larga_comparacion}\\
    \toprule
    \rowcolor{gobmxDorado} \encabezadodorado{Concepto} & \encabezadodorado{Valor 1} & \encabezadodorado{Valor 2} & \encabezadodorado{Nota} \\
    \midrule
    \endfirsthead

    \rowcolor{gobmxDorado} \encabezadodorado{Concepto} & \encabezadodorado{Valor 1} & \encabezadodorado{Valor 2} & \encabezadodorado{Nota} \\
    \midrule
    \endhead

    \bottomrule
    \endlastfoot

    \textbf{Fila 1} & 100 & 200 & Ejemplo \\
    \textbf{Fila 2} & 150 & 250 & Prueba \\
    \textbf{Fila 3} & 200 & 300 & Test \\
  \end{longtable}
\end{tabladoradoLargo}

\end{document}